%% CAMERA-READY -- WoAIS2 @ Middleware 2026.
%% CFP: ACM SIGPLAN style, 10pt; accepted paper, now limited to 6 pages
%% total (technical content + references) per organizer confirmation.
\documentclass[sigplan,10pt,natbib=false]{acmart}

\AtBeginDocument{%
  \providecommand\BibTeX{{Bib\TeX}}}

%% Rights management -- confirmed values from the ACM rights form.
\copyrightyear{2026}
\acmYear{2026}
\setcopyright{cc}
\setcctype{by}
\acmConference[WoAIS '26]{Workshop on AI and Serverless }{December 14--18, 2026}{Tarragona, Spain}
\acmBooktitle{Workshop on AI and Serverless (WoAIS '26), December 14--18, 2026, Tarragona, Spain}
\acmDOI{10.1145/3861043.3861046}
\acmISBN{979-8-4007-3108-2/2026/12}

\RequirePackage[
  datamodel=acmdatamodel,
  style=acmnumeric,
  ]{biblatex}
\addbibresource{refs.bib}

%% /proc paths are long and \texttt{} content is never hyphenated, which
%% overfills the narrow sigplan columns. \procpath permits a line break
%% after each slash; no hyphen is inserted, as is conventional for paths
%% and URLs (a trailing hyphen would read as part of the path).
\newcommand{\procpath}[1]{\texttt{/proc/\allowbreak<pid>/\allowbreak#1}}

\begin{document}

%% ---------------------------------------------------------------------
%% TITLE -- placeholder, revisit once the paper's scope is fully locked.
%% ---------------------------------------------------------------------
\title{How Much Do Co-Located Containers Really Share? A Pairwise
  Measurement Study}
\renewcommand{\shorttitle}{How Much Do Co-Located Containers Really Share?}

%% ---------------------------------------------------------------------
%% AUTHORS -- adjust affiliation details / ORCID as needed.
%% ---------------------------------------------------------------------
\author{Hanlin Wu}
\affiliation{%
  \institution{KU Leuven}
  \city{Leuven}
  \country{Belgium}}
\email{hanlin.wu@student.kuleuven.be}

\author{Eddy Truyen}
\affiliation{%
  \institution{KU Leuven, DistriNet}
  \city{Leuven}
  \country{Belgium}}
\email{eddy.truyen@kuleuven.be}

\author{Wouter Joosen}
\affiliation{%
  \institution{KU Leuven, DistriNet}
  \city{Leuven}
  \country{Belgium}}
\email{wouter.joosen@kuleuven.be} %% TODO: verify -- guessed from the standard KU Leuven address pattern.

\renewcommand{\shortauthors}{Wu, Truyen, and Joosen}

%% ---------------------------------------------------------------------
%% ABSTRACT
%% ---------------------------------------------------------------------
\begin{abstract}
This paper investigates for contemporary container platforms how much
memory saving can be attained by co-locating instances of one image on
the same node -- as a FaaS platform does to keep a warm pool ready, or
as a blue-green deployment does across consecutive versions, pinning
different tenants to the version they were validated against. The
expectation is that images
built from overlapping layers let the kernel back them with a single
physical copy.  That expectation has never been verified for a specific pair of
instances, and standard metrics cannot verify it: Proportional Set
Size (PSS) divides a page's cost by its mapcount, and so reports
sharing with \emph{any} process on the host rather than with the
neighbour a scheduler is reasoning about.
We measure it directly, intersecting the resident page frames of two named containers through
the kernel's \texttt{pagemap} interface. Between two consecutive
builds of a multi-tenant Java-based application, whose application layers are
byte-identical, the instances share 2.4\,MiB of mapped memory -- under
half a percent of either one's resident footprint. A static analysis
of roughly 100 releases each of two widely-used Java applications finds
that most consecutive release pairs share no image
snapshot whatsoever. The cause is structural rather than semantic. A
container runtime names each unpacked layer by the hash chain of every
layer beneath it, so a metadata-only difference early in the stack
gives byte-identical files above it fresh inodes, and the kernel
cannot deduplicate them. On current mainstream container platforms, a
sharing discount should not be assumed from image-layer similarity
alone, particularly once co-located versions are built independently.
Existing proposals that key image storage on content rather than
layer position deserve consideration as a way out.
\end{abstract}

%% ---------------------------------------------------------------------
%% CCS CONCEPTS -- regenerate the numeric IDs via
%% https://dl.acm.org/ccs.cfm before submission; the concept_desc text
%% below is a best-effort placeholder.
%% ---------------------------------------------------------------------
\begin{CCSXML}
<ccs2012>
 <concept>
  <concept_id>00000000.0000000.0000000</concept_id>
  <concept_desc>Computer systems organization~Cloud computing</concept_desc>
  <concept_significance>500</concept_significance>
 </concept>
 <concept>
  <concept_id>00000000.00000000.00000000</concept_id>
  <concept_desc>Software and its engineering~Software deployment</concept_desc>
  <concept_significance>400</concept_significance>
 </concept>
 <concept>
  <concept_id>00000000.00000000.00000000</concept_id>
  <concept_desc>Software and its engineering~Virtual machines</concept_desc>
  <concept_significance>300</concept_significance>
 </concept>
 <concept>
  <concept_id>00000000.00000000.00000000</concept_id>
  <concept_desc>Computer systems organization~Cloud computing</concept_desc>
  <concept_significance>200</concept_significance>
 </concept>
</ccs2012>
\end{CCSXML}

\ccsdesc[500]{Computer systems organization~Cloud computing}
\ccsdesc[400]{Software and its engineering~Software deployment}
\ccsdesc[300]{Software and its engineering~Virtual machines}

\keywords{container images, memory sharing, page cache, cgroups,
  multi-tenancy, serverless cost-effectiveness, Kubernetes, JVM}

\maketitle

%% =======================================================================
\section{Introduction}
\label{sec:intro}

Container-based platforms pack many containers onto one host to
improve resource utilization. FaaS platforms keep pools of warm
containers ready to absorb invocations without a cold-start penalty,
and co-locating instances of the same image reduces cold starts
further \cite{beni_coldstarts}. Blue-green deployments similarly keep
several versions of one service live, pinning each tenant to the
version they were validated against rather than forcing everyone onto
whatever is current; a placement policy can co-locate different
versions of one image on a node specifically to maximize this resource
sharing \cite{truyen2020flexible}.

A recurring proposal for increasing memory utilization is to credit
the memory co-located instances are expected to share: if two
containers are built from overlapping image layers, the kernel's page
cache can map the same physical pages into both address spaces instead
of duplicating them, lowering an added instance's marginal cost below
its declared footprint. Bravo Ferreira et al.\ state this
explicitly for FaaS workloads and propose a scheduler that maximizes
this assumed overlap \cite{bravo2017more}; Truyen et al.\ likewise
build it into a cost model for co-located versions, treating it as
substantial \cite{truyen2020flexible}.

Both rely, however, on the \procpath{smaps} interface or PSS, which
marks a page \emph{Shared} whenever more than one process on the host
maps it -- real but not \emph{attributable} information, since it
never names which neighbor a page is shared with.

Attribution is only half the problem: crediting shared layers also
assumes they resolve to common physical memory, which requires both
containers to map the same inode. Container runtimes name each
unpacked layer by the hash chain beneath it, so a difference anywhere
earlier in the chain -- even metadata only -- gives every later layer
a fresh inode, faulting identical files in twice. How often this
happens across a real application's releases, and how much sharing it
removes, has not been measured.

In this paper, we introduce a simple measurement technique that
intersects the resident physical page frames of two \emph{named}
containers directly, via \procpath{pagemap}, giving a pairwise,
ground-truth answer to how much two specific instances share right
now, independent of anything else running on the host. We apply it to
a real, evolving JVM-based service to measure how much two instances of
the identical image share while running, and how much of that survives
once the images differ -- as between two tenant-pinned versions of one
service -- complemented by a static characterization of image-layer
overlap across many consecutive versions of two applications.

%% =======================================================================
\section{Problem Statement}
\label{sec:background}

When a platform keeps several versions of one container image resident
on a node -- during a blue-green migration, or to hold tenants on the
version they were validated against -- the question that decides
capacity is whether the second instance costs a full memory footprint
or a discounted one. Existing cost and placement models answer with a
per-pair term: the fraction $f$ of libraries two arbitrary containers
are assumed to share \cite{bravo2017more}, or the \emph{Shared}
quantity by which two co-located versions of one service are
discounted \cite{truyen2020flexible}. Both are properties of a
\emph{specific pair of instances}, and both are consumed as model
input rather than measured. Establishing what that term is actually
worth is the problem this paper addresses.

The quantity has a precise definition. Linux shares identical content
between processes through its virtual-memory subsystem: each address
space is a set of virtual memory areas (VMAs) mapping anonymous
regions or file contents, and the frames backing a file-backed VMA
come from the kernel's file cache, indexed by the file's inode
\cite{mauerer2008linux}. Two processes share a \emph{file-backed}
page's frame only when both map content resolving to the same inode
\cite{inode}; bytes read from a separate copy, at a
different path and hence a different inode, are faulted in twice.
Residency in both processes is necessary but not sufficient. A shared
library is mapped \texttt{MAP\_PRIVATE}, so a write by either side
quietly replaces that side's copy of the page and leaves the other
mapping the original -- the page stays resident in both, but is no
longer one frame. What two instances share is therefore the set of
physical frames they \emph{both} hold resident, and that is what we
measure. Content a process reads without mapping falls outside the
definition, since the page cache deduplicates it when the inode
matches but it never enters a page table; our figures accordingly
bound mapped sharing rather than all sharing.

No existing interface reports this pairwise quantity. RSS counts every page a process
maps, whatever else maps it; PSS divides each page's cost by its
mapcount \cite{corbet2007pss}; \texttt{smaps} splits resident memory
into \emph{Shared} (mapcount $>1$ anywhere on the host) and
\emph{Private} bytes. All three capture the \emph{degree} of sharing,
not its \emph{identity}: mapcount records how many processes map a
page, never which ones. On a machine running only the pair of
interest the distinction collapses, unremarked; on a production node
also running a database, sidecars, and unrelated tenants, an
instance's Shared bytes include pages shared with none of the peers
under consideration, overstating what moving a peer away would
reclaim. Measuring whole hosts does not recover it either: one
configuration's realized saving does not decompose into the per-pair
quantities a cost model needs, and configurations are exponential in
instance count. What is needed is the physical page-frame numbers
(PFNs) both instances hold resident -- the quantity our technique
computes directly.

A cheaper proxy might seem possible: \procpath{maps} prints a
device and inode per mapping. But those columns are unreliable -- on
an overlay mount we observed three containers report three different
device numbers for one file -- and inode equality is necessary, not
sufficient: it says nothing about how many pages were faulted in, nor
does a shared inode survive one side writing to the page. Only
resident page frames settle how much two instances actually share.

%% =======================================================================
\section{Related Work}
\label{sec:related-work}

Bravo Ferreira et al.\ \cite{bravo2017more} first demonstrated that
shared image layers can produce genuine physical sharing, showing via
\texttt{pmap} a kernel-tracked \emph{Shared} byte count for two
containers union-mounted over a common library layer, none when the
same content sits in a private file system. They build a bin-packing
model around a free parameter $f$, the assumed fraction of shared
libraries, and show even a small $f$ yields most of the achievable
packing-density gain -- but never measure $f$ for a real application.
Truyen et al.\ \cite{truyen2020flexible} apply the same mechanism to
two \emph{consecutive versions} of one application resident on a node
during a migration step, reproducing that \texttt{smaps}-based
measurement for a cost model where demand is proportional to
co-located instances, discounted by assumed sharing.

Cheng et al.\ \cite{cheng2024slo} take layer-wise sharing as the basis
for a placement algorithm rather than a cost model. Measuring host
buffer/cache footprint as containers are spawned together or spread
across machines, they find it stays flat rather than growing with
instance count, and build heuristics that co-locate workflow functions
to exploit that effect. Their quantity, measured at the host level,
includes content a process reads without mapping, complementary to
ours. The workflows they place share a common runtime layer at the
bottom of each image, where shared-prefix snapshots are most likely to
survive. Our results say what happens once that prefix is not intact,
as when images are built independently: the benefit such an algorithm
assumes is more fragile than a single shared base image makes it look.

Medes \cite{saxena2022medes} also attributes redundancy to a specific
pair of containers, but differently: it fingerprints memory chunks by
content hash and merges matches regardless of whether the kernel's
inode-based cache would share them. Across this line of work, sharing
is confirmed possible, engineered, or built into a cost model, but
never measured for a real, evolving application, nor tied to the
mechanism that grants or withholds it. This paper measures that
quantity -- the fraction Bravo Ferreira et al.'s model calls $f$, and
Truyen et al.'s \emph{Shared} -- directly and attributably, for both
co-location patterns a scheduler must reason about.

%% =======================================================================
\section{Measurement Methodology}
\label{sec:methodology}

\subsection{Measurement Tools}
\label{sec:tool}

\paragraph{Pairwise PFN intersection.}
To measure realized sharing between any two specific, named
containers, we built a tool that, for each container, walks
\procpath{maps} to enumerate its virtual memory areas and reads
\procpath{pagemap} \cite{corbet2007pss} to resolve every resident
virtual page to the PFN backing it. Given two containers, the tool
computes the intersection of their resident PFN sets: any PFN present
in both sets is a physical frame genuinely mapped by both containers
at the moment of measurement, independent of whether the kernel's own
\texttt{smaps} bookkeeping happens to also attribute that page to some
other, unrelated mapper elsewhere on the host. It reports the
intersection broken down by mapped file, so sharing can be attributed
to the specific library or archive responsible, and separately
intersects each container's \emph{anonymous} memory, non-zero only if
same-page merging (KSM) is active -- ruling KSM in or out as a
confound rather than assuming it away. Reading \texttt{pagemap}'s
physical frame numbers requires elevated privileges, so the tool runs
on the node, not inside either measured container; we use it as an
offline, ground-truth instrument for testing cost-model assumptions,
not as a mechanism intended to run inside a live scheduler. Applying
it is independent of application or runtime, except for the
application-specific details below.

\paragraph{Static snapshot overlap.}
Whether two instances can share a mapped file at all is settled before
either runs, by how their images were built. Images are stacked as
layers, and when a container starts the runtime decompresses each
layer once into its own \emph{snapshot} directory and stitches the
relevant snapshots together through a union (overlay) file system
\cite{ociimagespec,kerneloverlay}. Overlay mounts route a
mapped file through to the underlying lower-layer file, so the inode
backing a mapped library is the one inside that snapshot directory.
Crucially, matching content is not enough to put two images in the
same snapshot: the runtime names each snapshot by its \emph{ChainID},
the cumulative hash of the uncompressed digests (diffIDs, from the
image configuration rather than the manifest) of every layer at or
below that position \cite{ociimagespec}. A shared snapshot at position
$k$ therefore requires layers $0,\dots,k$ to be identical. If any
earlier layer differs at all, every later layer is unpacked into a
fresh, separately-inoded snapshot regardless of its own content, so
two images can share every byte of an application's binaries and still
map them through different inodes. A copy-up reaches the same outcome
more simply: a container that writes to a file, or merely changes its
mode at entrypoint, gets a private copy with its own inode even when
the ChainIDs match. Both are properties of per-layer snapshotters such
as overlayfs under containerd and Docker's \texttt{overlay2};
content-addressed designs keying storage on file content instead,
such as composefs \cite{composefs}, can share identical files across
different layer chains.

Comparing the ChainIDs of two images from position 0 finds the first
position at which they diverge, and every layer before it would share
a snapshot were the two ever unpacked on one node. We define
\emph{snapshot overlap} between two images as the fraction of total
image bytes, weighted by the manifest's per-layer sizes, falling
inside that shared prefix: an upper bound, in advance of any runtime
measurement, on how much two images could possibly share.

To see how snapshot overlap depends on the size of an upgrade step, we
bucket consecutive-release pairs by the level of the semantic-version
identifier that changed between them. A \emph{build} step re-packages
the same (major, minor, patch) triple with a new build number (e.g.\
\texttt{26.1.2-1} $\to$ \texttt{26.1.2-2}); a \emph{patch} step changes
the last version component (\texttt{26.1.2} $\to$ \texttt{26.1.3}); a
\emph{minor} step changes the middle component (\texttt{26.1}
$\to$ \texttt{26.2}); and a \emph{major} step changes the first
(\texttt{25} $\to$ \texttt{26}). Within each granularity we take one
representative tag per version unit -- the newest build of each patch,
the newest patch of each minor, and so on -- and compute snapshot
overlap over every consecutive forward pair of those representatives.
The PFN-intersection tool, the release-tag overlap script, and the raw
measurements behind our figures are available at
\url{https://github.com/k8-scalar/memory-sharing-between-container-pairs}.

\subsection{Application and Workload}
\label{sec:app-workload}

We study Keycloak, an open-source identity and access management (IAM)
application, for three reasons: it is JVM-based, mapping shared native
libraries and a class-data-sharing archive; it is natively multi-tenant, its
\emph{Realm} construct
isolating each tenant in one instance, exactly what a blue-green
migration must reason about; and it has an actively maintained history
spanning dozens of versions, needed to test the mechanism across real,
independently built releases. To check that the layer-divergence
mechanism is not an artifact of one application, we extend the
static-overlap analysis to a second, independently maintained JVM
service, Elasticsearch. We collected roughly 100 recent release tags
each of Keycloak (\texttt{v20.x.x-x} through \texttt{v26.x.x-x}, 96
tags) and Elasticsearch (\texttt{v8.x.x} through \texttt{v9.x.x}, 121
tags); the
build granularity applies only to Keycloak, whose official image tags
carry an explicit build suffix. Keycloak
is not the highest-stakes case; that distinction belongs to serverless
model-serving, where mapped weights dwarf a
JVM's artifacts. This makes Keycloak a conservative choice: if the
mechanism collapses sharing here, it is unlikely to spare an
application with more at stake. We drive it with a fixed
per-tenant workload: a login request, several backend introspection
requests, a token refresh, more introspection requests, and a logout,
repeated by each simulated user for the duration of a run, with
randomized think-time between requests. Fixing this behavior keeps
memory measurements comparable across runs and versions. We report
results from a Kubernetes-hosted deployment of this workload, not from
a commercial FaaS platform's own execution environment.

\subsection{Deployment Topology}
\label{sec:topology}

The static-overlap script needs no cluster: it computes its bound
directly from image manifests, without running anything. The
PFN-intersection tool does, to measure what two instances actually
hold resident on a live deployment. Experiments ran on a three-node Kubernetes cluster (one control-plane,
two worker nodes), with all Keycloak instances and their PostgreSQL
database on one worker node and load-generation and monitoring on the
other. Workload was generated with K6 \cite{k6}, one pod per simulated
tenant running the request sequence described above. Pod-level CPU and
memory metrics were collected every 15\,s from each node's
\texttt{/proc} and \texttt{/sys} hierarchies.

%% =======================================================================
\section{Quantifying Memory Sharing}
\label{sec:results}

\subsection{Static Layer Overlap}
\label{sec:static-overlap}

Applying the static-overlap method of \S\ref{sec:tool} to Keycloak and
Elasticsearch gives Table~\ref{tab:static-overlap}.

\begin{table}[t]
  \centering
  \caption{Static snapshot overlap by version-step granularity: mean
    and max snapshot overlap as a percentage of image bytes, and the
    number of consecutive-release pairs $n$ each figure is computed
    over, and \emph{zero}, the number of those pairs sharing no
    snapshot at all. The build granularity applies only to Keycloak.}
  \label{tab:static-overlap}
  \begin{tabular}{lcccc}
    \toprule
     & Build & Patch & Minor & Major \\
    \midrule
    \multicolumn{5}{l}{\textbf{Keycloak} (v20--v26, 96 tags)} \\
    mean & 1.15\% & 1.34\% & 0.00\% & 0.00\% \\
    max  & 13.1\% & 13.1\% & 0.0\%  & 0.0\%  \\
    $n$  & 19     & 64     & 6      & 6      \\
    zero & 15     & 38     & 6      & 6      \\
    \midrule
    \multicolumn{5}{l}{\textbf{Elasticsearch} (v8--v9, 121 tags)} \\
    mean & \textemdash & 2.20\% & 0.87\% & 0.00\% \\
    max  & \textemdash & 5.7\%  & 5.6\%  & 0.0\%  \\
    $n$  & \textemdash & 97     & 22     & 1      \\
    zero & \textemdash & 49    & 18     & 1      \\
    \bottomrule
  \end{tabular}
\end{table}

The result is striking, and the means understate it: most consecutive
release pairs share no snapshot whatsoever (15 of 19 Keycloak build
pairs, 38 of 64 patch pairs, 49 of 97 Elasticsearch patch pairs, and
every Keycloak minor and major pair). The non-zero means come from a
minority retaining a small base-OS prefix, not sharing spread evenly
across the history. This bounds, ahead of any runtime measurement, how
much two consecutive releases could possibly share -- including the
multi-version tenancy case, where co-located versions are typically at
least a patch apart: runtime sharing below can therefore never exceed
a few percent once images are more than one patch apart, however
favorable the runtime conditions.

\subsection{Runtime Sharing: Same-Image Case}
\label{sec:same-image}

We co-located two instances of the identical Keycloak image on one
node across ten recent patch versions (v26.0.0 to v26.2.2), let each
pair settle for 30\,s, and measured PFN intersection between them.
Realized sharing is small and stable across versions: mean 25.8\,MiB
(s.d.\ 0.21\,MiB), around 3\% of the two-instance RSS union (RSS and
PSS both near 900\,MiB) and around 5\% of one instance's own RSS. We
report against RSS because RSS, not PSS, is what a cgroup memory limit
enforces; PSS's discount is a host-wide
average over however many processes map a page, not an amount
attributable to one co-located instance -- exactly the ambiguity our
pairwise technique replaces. Here the choice barely matters: sharing is
small enough that the RSS-sum and PSS-sum land within a percent of each
other.

Broken down by mapped file, shared bytes concentrate in a handful of
JVM-internal artifacts rather than application code, which matters for
how far the result generalizes: \texttt{libjvm.so} contributes roughly
16\,MiB (92\% of its own 17.4\,MiB footprint, by far the largest
share), \texttt{classes.jsa} \textasciitilde2.7\,MiB of 13.6\,MiB, and
the runtime module image \textasciitilde3.2\,MiB of 148\,MiB.
Keycloak's own \textasciitilde144\,MiB of application JARs contribute
\emph{nothing} to this figure -- not because they fail to share, but
because the class loader \texttt{pread}s rather than maps them, so
they never acquire a VMA to intersect. They are shared nonetheless:
\texttt{fincore} over the containers' snapshot directories shows the
316 JARs under \texttt{lib/main} resident in the page cache and held
once for both instances, some 122\,MiB invisible to a page-frame
intersection. The 25.8\,MiB of mapped sharing is thus a lower bound.
That 122\,MiB is also a more fragile form of sharing than mapped
sharing: the class loader reads each JAR once at startup and never
touches those pages again, so they sit in the page cache as
reclaimable, unmapped data the kernel can evict under memory pressure
and must re-read from disk on the next restart -- a one-time cost
recovered per boot, not a standing credit like a resident mapping.
Anonymous-page PFN intersection -- non-zero only if same-page merging
(KSM) were active -- was negligible on our nodes, ruling out KSM: the
sharing we observe is entirely file-cache sharing of mapped files, not
kernel deduplication of Keycloak's own heap.

\subsection{Runtime Sharing: Multi-Version Tenancy Case}
\label{sec:cross-version}

If realized sharing survives the smallest possible image change, it is
worth modeling in blue-green deployments. To find out, we compared two consecutive
\emph{builds} of the same Keycloak patch version, \texttt{v26.1.2-1}
and \texttt{v26.1.2-2}, which bundle the same OpenJDK build and are
about as close to identical as two independently tagged, independently
built releases ever get. We measured three pairs under a synchronized
400-VU load ramp: two same-build baselines, as a sanity-check control
against \S\ref{sec:same-image} rather than a result in their own right
(both instances at \texttt{-1}: 25.84\,MiB shared; both at
\texttt{-2}: 25.86\,MiB), and the cross-build pair, at 3\% static
overlap (Table~\ref{tab:static-overlap}'s build column): 2.38\,MiB
shared. The controls reproduce the same-image finding of
\S\ref{sec:same-image} almost exactly -- two independent measurements
of the identical configuration agree to within 0.02\,MiB. The
cross-build pair is the informative one. Both images bundle the same
OpenJDK build, so the two processes'
\procpath{maps} list the same 41 file paths -- there is no
sense in which the cross-build pair maps \emph{less} of the same
content than the same-build pairs do. Yet realized sharing drops by an
order of magnitude, from around 26\,MiB to 2.4\,MiB, and the 2.4\,MiB
that remains comes entirely from base-OS files
(\texttt{libc.so.6}, \texttt{libm.so.6}, \texttt{ld-linux}, locale
data) in the layer prefix the two images do still share.

To see exactly what changed, we unpacked all five image layers of both
builds and diffed every tar member. The two base layers (L0, L1) are
identical. Only seven files, 14\,MiB in total, differ in
\emph{content}, and all seven sit in L2, the distribution rootfs,
alongside 936 members differing only in \texttt{mtime}. Above it, L3 --
the layer that installs Keycloak itself, including every file whose
sharing we measured -- has all 444 of its regular files bit-for-bit
identical between the builds, but records six of its own
\texttt{mtime} differences, all on directories; L4 records three of its
own. Two causes compound here, each independently sufficient: L2's
content drift forces L3 and L4 into fresh, separately-inoded snapshots
via the ChainID cascade \emph{regardless of their own content}, and
L3/L4's own \texttt{mtime}-only differences would break their sharing
even if L2 matched. Pinning either layer's timestamps alone would not
restore sharing, since the other cause remains. We confirmed this
directly rather than inferring it: resolving each
container's overlay lowerdir chain and calling \texttt{stat} on the
resulting snapshot paths puts the two builds' copies of
\texttt{libjvm.so} at distinct inodes, versus one shared inode for two
instances of the same build, the files being byte-identical at
27{,}258{,}184 bytes either way. The unmapped content divides the same
way: the two builds' JAR sets resolve to different inodes, turning the
122\,MiB cached once for a same-image pair into two copies (122 and
131\,MiB). A manifest still showing large declared overlap -- L0's
22\,MiB base layer untouched, L3's 156\,MiB Keycloak-install layer
byte-identical -- thus yields almost none of the physical sharing it
suggests, purely because of where in the stack these differences sit.

%% =======================================================================
\section{Implications for Cost-effective Scheduling Policies}
\label{sec:implications}

These results argue against crediting assumed sharing in placement or
capacity decisions. Even the most favorable case maps only a few
percent of either instance's resident memory in common, and that
ceiling collapses by an order of magnitude once co-located versions
diverge, routine for multi-tenant platforms. A scheduler built on an
overlap-aware placement model \cite{bravo2017more} optimizes
for a largely absent benefit. The collapse is not purely a timestamp
artifact: L2 itself carries seven files of genuine content drift
between builds (\S\ref{sec:cross-version}), so pinning layer timestamps
alone would not recover the gap. Keying snapshot identity on
\emph{content} rather than layer \emph{position} -- composefs
\cite{composefs} being the most developed proposal -- addresses the
cause directly, recovering sharing up to the same-image ceiling without
demanding reproducible builds from every producer.

Structurally, crediting sharing inside an instance's declared memory
limit is unsafe however much sharing exists. A PSS-style credit
assumes a shared page's cost splits between its mappers, but cgroups
charge a memory area in full to whichever cgroup instantiated it; for
an area used from different cgroups, \emph{``to which cgroup the area
will be charged is in-deterministic''} \cite{kernelcg}. A limit set
below an instance's full footprint, assuming part of a shared region is
billed to a neighbour, can fail either way: the kernel may OOM-kill the
instance charged the whole region while the other runs with unbudgeted
headroom -- consistent with Kubernetes requiring request and limit to
be equal under its strictest QoS class \cite{k8sqos}.

Although out of scope for our core PFN measurement, we also examined
this channel: \S\ref{sec:same-image} measured 122\,MiB of JAR
page-cache content shared for a same-image pair, and
\S\ref{sec:cross-version} showed it collapses across versions for the
same inode-divergence reason as mapped sharing. Recent work builds a
placement algorithm on exactly this channel: Cheng et al.\
\cite{cheng2024slo} report that the host buffer/cache footprint of
five co-located containers of one function stays flat instead of
growing fivefold, and place workflow functions to maximize that
effect, measured at the host level rather than attributed to a
specific pair. Neither finding rescues the cost-effectiveness argument: being clean,
unwritten data, it may be reclaimed under pressure rather than
committed, and what survives still collapses across versions like
mapped sharing. A Kubernetes request or OpenWhisk slot, fixed in
advance, accounts for none of it, and saving it lowers host pressure
without admitting another container.

The stakes rise sharply for serverless inference, where a model
server's resident memory is dominated not by a language runtime but by
model weights: hundreds of megabytes to tens of gigabytes and
byte-identical across every instance serving it. Zero-copy
formats such as safetensors, or the GGUF files \texttt{llama.cpp}
maps, enter the address space through \texttt{mmap} and fall inside
the quantity measured here; loaders that read weights into anonymous
memory land in the page-cache channel instead. Either way the inode
condition decides whether two co-located instances pay once or twice,
so the mechanism of \S\ref{sec:static-overlap} applies unchanged --
over objects three orders of magnitude larger, often delivered through
mounted volumes rather than image layers, changing which rule governs
the inode. That holds only for weights resident in host RAM, though:
once copied into GPU memory, each instance holds an independent device
allocation outside the host page cache, a separate problem this
mechanism cannot reach. Measuring the host-RAM case is next.

%% =======================================================================
\section{Threats to Validity}
\label{sec:threats}

Our measurement counts only pages both instances map; unmapped content
is still deduplicated when the inode matches
(\S\ref{sec:same-image}), a channel reported as substantial for
small function images \cite{cheng2024slo}. Our figures are therefore a
lower bound, smaller for a runtime that maps most of what it reads and
larger for one that does not. The measurements are strictly pairwise; host-level
saving across more than two instances is bounded by one cached copy.
The page-cache figures of \S\ref{sec:same-image}--\ref{sec:cross-version}
come from one exemplary set on an idle node; characterizing that
channel under memory pressure is future work.

Our measurements come from one application family, a JVM service whose
shareable memory sits in a few memory-mapped native libraries and a
class-data-sharing archive. A different runtime profile -- a
statically-linked binary, or one mapping more of its own files -- could
realize a different fraction; the mechanism is not JVM-specific, but
the magnitude is. The cross-version experiment rests on three pairs,
adequate given the 0.02\,MiB spread between same-build controls. Results
come from Kubernetes, not a commercial FaaS environment; generalization
rests on the kernel and OCI mechanism, both orchestrator-independent.

%% =======================================================================
\section{Conclusion}
\label{sec:conclusion}

Crediting assumed memory sharing between co-located containers raises
resource utilization, but the credit has never been measured precisely
for a specific pair. Using pairwise page-frame intersection, we
measured it for a JVM-based service in two situations: concurrent
instances of one version, and tenant-pinned versions on a host. Sharing
is small in the first case and collapses by an order of magnitude in
the second, driven by how snapshots are identified, not by the
application. A sharing discount should not be assumed from
layer similarity alone, especially once co-located versions are
built independently.

%% ---------------------------------------------------------------------
%% GENERATIVE AI DISCLOSURE.
%% ---------------------------------------------------------------------
\begin{acks}
This research is partially funded by the Internal Funds KU Leuven. An
AI assistant (Claude, Anthropic) was used while preparing this
manuscript: for discussion and verification of the measurement
methodology and of related work, for drafting and revising prose under
the authors' direction, and for assembling the \LaTeX{} scaffold. All
experimental results, data, and conclusions are the authors' own, and
AI-assisted text was reviewed and edited by the authors, who take full
responsibility for the submitted content.
\end{acks}

{\footnotesize
\printbibliography
}

\end{document}
